% GENERATED FILE. Edit canonical lessons, then run npm run generate:books.
% canonical-source-hash: fnv1a64:6331d84757ca27fd
% canonical-lessons: KA-C127-samudra, KA-C127-kadu, KA-C127-mannu, KA-C127-kallu, KA-C127-gudugu

\chapter{Sea, Forest, Soil, Stone, Thunder}
\label{ch:ka-sea-forest-soil-stone-thunder}
\hlchaptermodality{127}

\begin{chapteropening}
\textbf{By the end of this chapter:} \emph{I can say the sea, a forest, soil, a stone and thunder.}

Complete the last lesson of chapter 127: i can say the sea, a forest, soil, a stone and thunder.
\end{chapteropening}

\section[samudra]{\kn{ಸಮುದ್ರ} (samudra) — the sea}

\label{lesson:KA-C127-samudra}



\begin{quote}
Before the new one: say the Kannada for pepper, chilli, then the Kannada for cumin.
\end{quote}

\subsection*{You'll want to know: \kn{ಸಮುದ್ರ}}
\textbf{\kn{ಸಮುದ್ರ}} — \emph{samudra} — \textquotedblleft{}the sea\textquotedblright{}.

\begin{grammarlens}[title={What you've built}]
One more word for this chapter.
\end{grammarlens}

\subsection*{Your turn}
\begin{itemize}
\raggedright
  \item \emph{Say it:} \emph{samudra}
  \item \emph{Say it:} \emph{samudra}, once more
  \item \emph{Say it:} say \emph{samudra}
  \item \emph{Recall:} say the Kannada for pepper, chilli, then the Kannada for cumin, then say \emph{samudra} again
\end{itemize}

\begin{tcolorbox}[breakable,colback=teal!4,colframe=teal!35!black,title={Before you move on}]
What does \kn{ಸಮುದ್ರ} mean? (\textquotedblleft{}The sea\textquotedblright{}.) Say it once more.
\end{tcolorbox}
\section[kāḍu]{\kn{ಕಾಡು} (kāḍu) — a forest}

\label{lesson:KA-C127-kadu}



\begin{quote}
Before the new one: say the Kannada for cumin, then the Kannada for the sea.
\end{quote}

\subsection*{You'll want to know: \kn{ಕಾಡು}}
\textbf{\kn{ಕಾಡು}} — \emph{kāḍu} — \textquotedblleft{}a forest\textquotedblright{}.

\begin{grammarlens}[title={What you've built}]
One more word for this chapter.
\end{grammarlens}

\subsection*{Your turn}
\begin{itemize}
\raggedright
  \item \emph{Say it:} \emph{kāḍu}
  \item \emph{Say it:} \emph{kāḍu}, once more
  \item \emph{Say it:} say \emph{kāḍu}
  \item \emph{Recall:} say the Kannada for cumin, then the Kannada for the sea, then say \emph{kāḍu} again
\end{itemize}

\begin{tcolorbox}[breakable,colback=teal!4,colframe=teal!35!black,title={Before you move on}]
What does \kn{ಕಾಡು} mean? (\textquotedblleft{}A forest\textquotedblright{}.) Say it once more.
\end{tcolorbox}
\section[maṇṇu]{\kn{ಮಣ್ಣು} (maṇṇu) — soil, earth}

\label{lesson:KA-C127-mannu}



\begin{quote}
Before the new one: say the Kannada for the sea, then the Kannada for a forest.
\end{quote}

\subsection*{You'll want to know: \kn{ಮಣ್ಣು}}
\textbf{\kn{ಮಣ್ಣು}} — \emph{maṇṇu} — \textquotedblleft{}soil, earth\textquotedblright{}.

\begin{grammarlens}[title={What you've built}]
One more word for this chapter.
\end{grammarlens}

\subsection*{Your turn}
\begin{itemize}
\raggedright
  \item \emph{Say it:} \emph{maṇṇu}
  \item \emph{Say it:} \emph{maṇṇu}, once more
  \item \emph{Say it:} say \emph{maṇṇu}
  \item \emph{Recall:} say the Kannada for the sea, then the Kannada for a forest, then say \emph{maṇṇu} again
\end{itemize}

\begin{tcolorbox}[breakable,colback=teal!4,colframe=teal!35!black,title={Before you move on}]
What does \kn{ಮಣ್ಣು} mean? (\textquotedblleft{}Soil, earth\textquotedblright{}.) Say it once more.
\end{tcolorbox}
\section[kallu]{\kn{ಕಲ್ಲು} (kallu) — a stone}

\label{lesson:KA-C127-kallu}



\begin{quote}
Before the new one: say the Kannada for a forest, then the Kannada for soil.
\end{quote}

\subsection*{You'll want to know: \kn{ಕಲ್ಲು}}
\textbf{\kn{ಕಲ್ಲು}} — \emph{kallu} — \textquotedblleft{}a stone\textquotedblright{}.

\begin{grammarlens}[title={What you've built}]
One more word for this chapter.
\end{grammarlens}

\subsection*{Your turn}
\begin{itemize}
\raggedright
  \item \emph{Say it:} \emph{kallu}
  \item \emph{Say it:} \emph{kallu}, once more
  \item \emph{Say it:} say \emph{kallu}
  \item \emph{Recall:} say the Kannada for a forest, then the Kannada for soil, then say \emph{kallu} again
\end{itemize}

\begin{tcolorbox}[breakable,colback=teal!4,colframe=teal!35!black,title={Before you move on}]
What does \kn{ಕಲ್ಲು} mean? (\textquotedblleft{}A stone\textquotedblright{}.) Say it once more.
\end{tcolorbox}
\section[guḍugu]{\kn{ಗುಡುಗು} (guḍugu) — thunder}

\label{lesson:KA-C127-gudugu}



\begin{quote}
Before the new one: say the Kannada for soil, then the Kannada for a stone.
\end{quote}

\subsection*{You'll want to know: \kn{ಗುಡುಗು}}
\textbf{\kn{ಗುಡುಗು}} — \emph{guḍugu} — \textquotedblleft{}thunder\textquotedblright{}.

\begin{grammarlens}[title={What you've built}]
One more word for this chapter.
\end{grammarlens}

\subsection*{Your turn}
\begin{itemize}
\raggedright
  \item \emph{Say it:} \emph{guḍugu}
  \item \emph{Say it:} \emph{guḍugu}, once more
  \item \emph{Say it:} say \emph{guḍugu}
  \item \emph{Recall:} say the Kannada for soil, then the Kannada for a stone, then say \emph{guḍugu} again
\end{itemize}

\begin{tcolorbox}[breakable,colback=teal!4,colframe=teal!35!black,title={Before you move on}]
What does \kn{ಗುಡುಗು} mean? (\textquotedblleft{}Thunder\textquotedblright{}.) Say it once more.
\end{tcolorbox}
